\documentclass[12pt,a4paper]{article}
\usepackage[utf8]{inputenc}
\usepackage{amsmath,amssymb,amsthm}
\usepackage{geometry}
\usepackage{hyperref}
\usepackage{algorithm}
\usepackage{algpseudocode}
\usepackage{graphicx}
\usepackage{mathtools}
\usepackage{tabularx}
\usepackage{booktabs}
\usepackage{caption}
\usepackage{cleveref}
\usepackage{longtable}
\newcolumntype{C}{>{\centering\arraybackslash}X}

\geometry{margin=1in}

\theoremstyle{plain}
\newtheorem{theorem}{Theorem}[section]
\newtheorem{lemma}[theorem]{Lemma}
\newtheorem{proposition}[theorem]{Proposition}
\newtheorem{corollary}[theorem]{Corollary}

\theoremstyle{definition}
\newtheorem{definition}[theorem]{Definition}
\newtheorem{example}[theorem]{Example}

\theoremstyle{remark}
\newtheorem{remark}[theorem]{Remark}

\DeclareMathOperator*{\argmin}{arg\,min}
\DeclareMathOperator{\sign}{sign}
\DeclareMathOperator{\diag}{diag}
\DeclareMathOperator{\prox}{prox}

\title{Sheaf-Theoretic Methods in Groundwater Hydrogeochemistry: \\
A Unified Mathematical Framework for Network Topology Inference, \\ Geochemical Inverse Modeling, and Isotope Source Apportionment}

\author{Dickson Abdul-Wahab}
\date{\today}

\begin{document}

\maketitle

% FRONT MATTER
\section*{Disclaimer}
The software and methods described in this document are provided "as is" without warranty of any kind. The authors and contributors assume no responsibility for errors or omissions in the software or documentation, or for damages resulting from the use of the information contained herein.

\begin{abstract}
	We present a comprehensive mathematical framework for integrated hydrogeochemical analysis in groundwater systems using sheaf-theoretic principles. The framework provides a unified computational approach for: (1) inferring flow network topology and connectivity from incomplete hydraulic head, topographic, and physics-based priors; (2) solving reaction path inverse problems to identify optimal transport processes and parsimonious geochemical reaction sets along network edges; and (3) performing trace-element forensics through Bayesian isotope mixing models and vadose zone travel-time deconvolution. By combining weighted least squares optimization with non-negativity constraints, sparse LASSO regression, thermodynamic constraints from PHREEQC, and Richards equation simulation, the method produces physically consistent, parsimonious models suitable for both reactive transport characterization and contamination source apportionment. We establish theoretical results on the optimality of transport solutions and the convergence of the underlying coordinate descent algorithms. Similar to other open-source projects, the implementation is available at \url{https://github.com/dabdul-wahab1988/Hydrosheaf}.
\end{abstract}

\newpage
\tableofcontents
\newpage
\listoffigures
\listoftables
\newpage

\section*{List of Main Variables}
\addcontentsline{toc}{section}{List of Main Variables}

\begin{longtable}{@{} l p{0.75\textwidth} @{}}
	\toprule
	\textbf{Symbol}                                           & \textbf{Description}                                                                                                         \\
	\midrule
	$G = (V, E)$                                              & Directed graph representing flow network (vertices = wells, edges = flow paths)                                              \\
	$\mathbf{x}_v \in \mathbb{R}^n$                           & Ion concentration vector at node $v$ (implementation default $n=10$: Ca, Mg, Na, HCO$_3$, Cl, SO$_4$, NO$_3$, F, Fe, PO$_4$) \\
	$\boldsymbol{\theta}$                                     & Transport model parameters ($\gamma$ for evaporation, $f$ for mixing)                                                        \\
	$\gamma \geq 1$                                           & Evaporation concentration factor                                                                                             \\
	$f \in [0,1]$                                             & Mixing fraction with endmember $\mathbf{x}_{\text{end}}$                                                                     \\
	$S \in \mathbb{R}^{n \times m}$                           & Stoichiometric matrix ($m$ reactions, $n$ ions)                                                                              \\
	$\mathbf{z} \in \mathbb{R}^m$                             & Reaction extent vector (positive = dissolution, negative = precipitation)                                                    \\
	$\mathbf{W} = \diag(w_1, \ldots, w_n)$                    & Diagonal weight matrix (typically inverse variance)                                                                          \\
	$\lambda > 0$                                             & L1 regularization parameter (controls sparsity)                                                                              \\
	$\boldsymbol{\ell}, \mathbf{u} \in \mathbb{R}^m$          & Lower and upper bounds on reaction extents (from saturation indices)                                                         \\
	$A(\boldsymbol{\theta}), \mathbf{b}(\boldsymbol{\theta})$ & Affine transformation representing transport process                                                                         \\
	$\mathbf{r} \in \mathbb{R}^n$                             & Residual vector after transport ($\mathbf{x}_v - A\mathbf{x}_u - \mathbf{b}$)                                                \\
	$\mathcal{P}_*$                                           & Penalty functions (isotope, EC/TDS, charge balance, thermodynamic)                                                           \\
	$\text{SI}$                                               & Saturation index $= \log_{10}(\text{IAP}/K_{sp})$                                                                            \\
	$h_i, h_j$                                                & Hydraulic heads at nodes $i, j$                                                                                              \\
	$p(i \to j)$                                              & Probability of flow from node $i$ to $j$                                                                                     \\
	$\pi_{uv}$                                                & Evaporation gate for sheaf topology refinement (0 = conservative, 1 = evaporative)                                           \\
	$\mathcal{E}_{\text{sheaf}}$                              & Sheaf consistency energy used for isotope-based topology refinement                                                          \\
	$\delta^{18}\text{O}, \delta^2\text{H}$                   & Stable water isotope ratios (permil relative to VSMOW)                                                                       \\
	$d = \delta^2\text{H} - 8 \cdot \delta^{18}\text{O}$      & Deuterium excess (permil)                                                                                                    \\
	$\mathcal{S}(\alpha, \tau)$                               & Soft-thresholding operator for LASSO                                                                                         \\
	$T$                                                       & Number of coordinate descent iterations to convergence                                                                       \\
	$|\mathcal{C}|$                                           & Number of candidate transport models                                                                                         \\
	$\mathcal{A} = \{j : z_j^* \neq 0\}$                      & Active set of reactions (non-zero extents)                                                                                   \\
	$R^2$                                                     & Coefficient of determination (fit quality metric)                                                                            \\
	\bottomrule
\end{longtable}

Throughout the document, vectors are denoted in boldface (e.g., $\mathbf{x}, \mathbf{z}$), scalars in italic (e.g., $\gamma, \lambda$), and matrices in boldface uppercase (e.g., $\mathbf{W}, S$). All concentration units are mmol/L unless otherwise specified.

\newpage

% SECTION 1
\section{Introduction}\label{sec:intro}

\subsection{Overview}\label{sec:overview}
Groundwater hydrogeochemistry seeks to understand the chemical evolution of water as it flows through aquifer systems. The inverse problem---determining which combination of transport processes (mixing, evaporation) and geochemical reactions (mineral dissolution/precipitation, redox reactions, ion exchange) explains observed chemical changes---is fundamental to aquifer characterization, contamination assessment, and water resource management.

Consider a groundwater flow network represented as a directed graph $G = (V, E)$, where vertices $V$ represent sampling locations (wells) and edges $E$ represent inferred flow connections. At each vertex $v \in V$, we observe a vector of ion concentrations $\mathbf{x}_v \in \mathbb{R}^n$ (default $n=10$ columns: Ca, Mg, Na, HCO$_3$, Cl, SO$_4$, NO$_3$, F, Fe, PO$_4$ in mmol/L).

For each directed edge $(u,v) \in E$, we seek to determine:
\begin{enumerate}
	\item The dominant transport process (evaporation or mixing with an endmember)
	\item The extent of mineral reactions that explain residual chemical changes
	\item The consistency of this explanation with thermodynamic, isotopic, and mass balance constraints
\end{enumerate}

To obtain meaningful, physically plausible solutions, we employ three complementary strategies:
\begin{itemize}
	\item \textbf{Sparsity regularization}: Prefer explanations involving fewer reactions (Occam's razor). The L1 penalty in LASSO regression automatically selects a sparse subset of reactions from the candidate dictionary.
	\item \textbf{Thermodynamic constraints}: Enforce mineral equilibrium bounds via saturation indices computed by PHREEQC. This eliminates geochemically impossible solutions.
	\item \textbf{Multi-physics integration}: Incorporate isotope data ($\delta^{18}$O, $\delta^2$H) and electrical conductivity to discriminate between processes.
\end{itemize}

\subsection{Comparison with Existing Tools}\label{sec:comparison}
The framework presented in this document differs from existing tools in several important ways:

\begin{itemize}
	\item \textbf{Global Parsimony (Occam's Razor)}: Unlike PHREEQC's combinatorial approach which yields all mathematically possible models (often hundreds), Hydrosheaf uses L1 regularization (LASSO) to automatically select the simplest, most physically plausible reaction set.
	\item \textbf{Automated Topology}: Existing codes require the user to explicitly define "Flow Path A to B." Hydrosheaf infers the flow network structure probabilistically from hydraulic heads and topography, making it suitable for regional-scale studies where flow paths are uncertain.
	\item \textbf{Dual-Tier Source Apportionment}: Hydrosheaf is uniquely designed for contaminant forensics, integrating a specialized Bayesian engine for nitrate source discrimination that seamlessly blends hydrochemistry with isotope data.
\end{itemize}

\begin{table}[htbp]
	\centering
	\caption{Comparison of Inverse Modeling Capabilities and Features}
	\label{tab:software_comparison}
	\begin{tabularx}{\textwidth}{@{} l C C C C @{}}
		\toprule
		\textbf{Feature}           & \textbf{Hydrosheaf}                & \textbf{PHREEQC}     & \textbf{NETPATH} & \textbf{GWB}        \\
		\midrule
		\textbf{Algorithm}         & \textit{Sparse L1 Optimization}    & Combinatorial Search & Mass Balance     & Mass Balance        \\ \addlinespace
		\textbf{Network Inference} & \textit{Probabilistic Graph}       & User-Defined 1D Path & User-Defined     & User-Defined        \\ \addlinespace
		\textbf{Transport Logic}   & \textit{Hypothesis Comp. (AIC)}    & Fixed Assumption     & Fixed Mixing     & Fixed               \\ \addlinespace
		\textbf{Isotopes}          & \textit{Dual-Isotope Mixing}       & Basic Fractionation  & Rayleigh         & Basic               \\ \addlinespace
		\textbf{Uncertainty}       & \textit{Bayesian MCMC / Bootstrap} & Sensitivity Analysis & None             & Monte Carlo (React) \\ \addlinespace
		\textbf{Nitrate Source}    & \textit{CoDA + Bayesian Gating}    & None                 & None             & None                \\
		\bottomrule
	\end{tabularx}
\end{table}

\subsection{System Architecture: Sheaf-Theoretic Perspective}\label{sec:sheaf}
The sheaf-theoretic viewpoint interprets the groundwater network as a cellular complex where local data (ion concentrations at wells) must satisfy global consistency conditions (mass balance, charge balance, thermodynamic equilibrium). While we do not develop full sheaf cohomology here, the computational framework enforces consistency through residual minimization along edges, analogous to minimizing the discrepancy in sheaf-theoretic data structures.

\subsection{Input-Output Contract}\label{sec:io_contract}
The framework takes as input:
\begin{itemize}
	\item Ion concentration vectors (default $n=11$ columns: Ca, Mg, Na, K, HCO3, Cl, SO4, NO3, F, Fe, PO4).
	\item Supporting bulk water fields: EC, TDS, and pH.
	\item Optional isotope data ($\delta^{18}$O, $\delta^2$H, $\delta^{15}$N--NO$_3$, $\delta^{18}$O--NO$_3$).
	\item Hydraulic head measurements or topographic elevations for network inference.
	\item Candidate mineral reactions and endmember compositions.
	\item Optional temporal time-series CSV when \texttt{--temporal-enabled} is used.
	\item Optional physics priors from mesh-based flow/particle tracking (FloPy/MODPATH).
\end{itemize}

The framework produces as output:
\begin{itemize}
	\item For each edge $(u,v) \in E$: the optimal transport model (evaporation factor $\gamma$ or mixing fraction $f$).
	\item A sparse set of active reactions with quantified extents $\mathbf{z}$.
	\item Model fit quality metrics ($R^2$, residual norms).
	\item Diagnostic flags (charge balance errors, thermodynamic violations).
	\item Optional temporal summaries and uncertainty intervals.
\end{itemize}

\newpage

% SECTION 2
\section{Governing Mathematical Framework}\label{sec:math_framework}

\subsection{Graph-Theoretic Foundation}\label{sec:graph_foundation}

\subsubsection{Probabilistic Edge Weights}
In groundwater networks, flow directions are inferred from hydraulic head measurements. Given hydraulic heads $h_i, h_j$ with uncertainties $\sigma_i, \sigma_j$ at nodes $i, j$ separated by distance $d_{ij}$, the head difference is $\Delta h_{ij} = h_i - h_j \sim \mathcal{N}(\mu_{ij}, \sigma_{ij}^2)$. The probability of flow from $i$ to $j$ is:
\begin{equation}\label{eq:edge_probability}
	p(i \to j) = \Phi(\mu_{ij} / \sigma_{ij})
\end{equation}

\subsubsection{Hydraulic Gradient Guard}
To avoid spurious connections, we apply a gradient threshold. If $g_{ij} = |\Delta h_{ij}|/d_{ij} < g_{\min}$, we attenuate the edge probability:
\begin{equation}\label{eq:gradient_guard}
	p(i \to j) \leftarrow 0.5 + (p(i \to j) - 0.5) \cdot \frac{g_{ij}}{g_{\min}}
\end{equation}

\subsubsection{Three-Dimensional Extensions}
For 3D systems, the effective hydrological distance $d_{3D}$ uses a vertical anisotropy factor $\alpha_v \approx \sqrt{K_v/K_h}$:
\begin{equation}
	d_{3D} = \sqrt{(x_1-x_2)^2 + (y_1-y_2)^2 + \frac{(z_1-z_2)^2}{\alpha_v^2}}
\end{equation}
Wells with overlapping screened intervals are assigned higher connectivity probabilities via a screen overlap bonus.

\subsubsection{Sheaf Consistency and the Laplacian}
In the sheaf-theoretic view, we define stalks $\mathcal{F}(v) = \mathbb{R}^n$ at each vertex and restriction maps $\varphi_{ev}: \mathcal{F}(v) \to \mathcal{F}(e)$ at each edge. For a flow path $u \to v$, the restriction maps incorporate the transport affine transformation: $\varphi_{eu}(\mathbf{x}_u) = A \mathbf{x}_u + \mathbf{b}$ and $\varphi_{ev}(\mathbf{x}_v) = \mathbf{x}_v - R\mathbf{z}$. A global assignment is a 0-section $\mathbf{X} \in \bigoplus \mathcal{F}(v)$. The sheaf consistency energy is defined by the quadratic form:
\begin{equation}\label{eq:sheaf_energy}
	\mathcal{E}(\mathbf{X}) = \sum_{e=(u,v) \in E} w_e \|\varphi_{ev}(\mathbf{x}_v) - \varphi_{eu}(\mathbf{x}_u)\|^2
\end{equation}
This energy is equivalent to $\mathbf{X}^\top L \mathbf{X}$, where $L$ is the \textbf{Sheaf Laplacian}. Minimizing $\mathcal{E}$ corresponds to finding a section that is most consistent with the underlying geochemical and transport physics.

\subsection{Transport Processes}\label{sec:transport_processes}

\begin{remark}[Conservative Weighting]
	The weight matrix $\mathbf{W}$ refers to the conservative weighting configuration ($\mathbf{W}_{\text{cons}}$), typically prioritizing non-reactive tracers (Cl, Br, isotopes).
\end{remark}

\subsubsection{Evaporation Model}
The evaporation model assumes conservative scaling of all solute concentrations:
\begin{equation}\label{eq:evaporation}
	\mathbf{x}' = \gamma \mathbf{x}_u, \quad \gamma \geq 1
\end{equation}
\begin{theorem}[Optimal Evaporation Factor]\label{thm:evap_optimal}
	For the weighted least squares problem $\min_{\gamma \geq 1} \|\mathbf{x}_v - \gamma \mathbf{x}_u\|_{\mathbf{W}}^2$, the optimal solution is:
	\begin{equation}\label{eq:gamma_optimal}
		\gamma^* = \max\left\{1, \frac{\mathbf{x}_u^\top \mathbf{W} \mathbf{x}_v}{\mathbf{x}_u^\top \mathbf{W} \mathbf{x}_u}\right\}
	\end{equation}
\end{theorem}

\subsubsection{Mixing Model}
The mixing model represents dilution or mixing with a distinct water type:
\begin{equation}\label{eq:mixing}
	\mathbf{x}' = (1-f) \mathbf{x}_u + f \mathbf{x}_{\text{end}}, \quad f \in [0, 1]
\end{equation}
\begin{theorem}[Optimal Mixing Fraction]\label{thm:mixing_optimal}
	For the weighted least squares problem
	\[\min_{f \in [0,1]} \|\mathbf{x}_v - (1-f)\mathbf{x}_u - f\mathbf{x}_{\text{end}}\|_{\mathbf{W}}^2,\]
	the optimal solution is:
	\begin{equation}\label{eq:f_optimal}
		f^* = \max\left\{0, \min\left\{1, \frac{\mathbf{d}^\top \mathbf{W} (\mathbf{x}_v - \mathbf{x}_u)}{\mathbf{d}^\top \mathbf{W} \mathbf{d}}\right\}\right\}
	\end{equation}
	where $\mathbf{d} = \mathbf{x}_{\text{end}} - \mathbf{x}_u$.
\end{theorem}

\subsection{Geochemical Reactions}\label{sec:geochemical_reactions}

\subsubsection{Sparse Reaction Formulation (LASSO)}
After transport, the residual mass $\mathbf{r} = \mathbf{x}_v - A(\boldsymbol{\theta}^*)\mathbf{x}_u - \mathbf{b}(\boldsymbol{\theta}^*)$ is explained by mineral reactions. The LASSO problem seeks sparse reaction extents $\mathbf{z} \in \mathbb{R}^m$:
\begin{equation}\label{eq:lasso}
	\min_{\mathbf{z}} \frac{1}{2} \|\mathbf{r} - S\mathbf{z}\|_{\mathbf{W}}^2 + \lambda \|\mathbf{z}\|_1 \quad \text{subject to} \quad \boldsymbol{\ell} \leq \mathbf{z} \leq \mathbf{u}
\end{equation}
The L1 penalty $\|\mathbf{z}\|_1$ promotes sparsity, yielding a parsimonious explanation involving few reactions.

\paragraph{Derivation of the Coordinate Update}
To solve \eqref{eq:lasso}, we consider the univariate subproblem for $z_j$, holding all other $z_{k \neq j}$ fixed. Let $\mathbf{r}_j = \mathbf{r} - \sum_{k \neq j} z_k \mathbf{s}_k$ be the partial residual. The subproblem is:
\begin{equation}
	\min_{z_j} \frac{1}{2} w_j (r_j - s_j z_j)^2 + \lambda |z_j|
\end{equation}
Setting the subdifferential to zero, $0 \in -s_j w_j (r_j - s_j z_j) + \lambda \partial |z_j|$, where $\partial |z_j|$ is the signum function if $z_j \neq 0$ and $[-1, 1]$ if $z_j = 0$. Solving for $z_j$ yields the \textbf{soft-thresholding operator}:
\begin{equation}\label{eq:soft_thresh}
	z_j = \frac{\mathcal{S}(s_j w_j r_j, \lambda)}{s_j^2 w_j}, \quad \mathcal{S}(\rho, \lambda) = \text{sign}(\rho) \max(0, |\rho| - \lambda)
\end{equation}
In the multi-ion case with diagonal weights $\mathbf{W}$, the term $s_j w_j r_j$ generalizes to the dot product $\mathbf{s}_j^\top \mathbf{W} \mathbf{r}_j$.

\subsubsection{Thermodynamic Constraints and Stoichiometry}
Saturation indices (SI) computed by PHREEQC determine thermodynamically feasible reactions bounds $(\ell_j, u_j)$:
\begin{equation}\label{eq:si_bounds}
	(\ell_j, u_j) = \begin{cases}
		(0, +\infty)       & \text{if } \text{SI}_u < -\tau \text{ and } \text{SI}_v < -\tau \quad \text{(dissolution only)} \\
		(-\infty, 0)       & \text{if } \text{SI}_v > \tau \quad \text{(precipitation only)}                                 \\
		(-\infty, +\infty) & \text{otherwise} \quad \text{(free)}
	\end{cases}
\end{equation}
The framework explicitly distinguishes between \textbf{Open System} (infinite $CO_2$ reservoir) and \textbf{Closed System} (limited $CO_2$) carbonate dissolution. For Calcite, the open stoichiometry is $CaCO_3 + CO_2 + H_2O \to Ca^{2+} + 2HCO_3^-$ (1:2 ratio), while the closed stoichiometry is $CaCO_3 + H^+ \to Ca^{2+} + HCO_3^-$ (1:1 ratio). Similar variants are provided for Dolomite and Magnesite.

Layer-specific mineral dictionaries are supported for 3D stratified systems.

\subsection{Isotope Hydrogeology}\label{sec:isotope_physics}

\subsubsection{Local Meteoric Water Line}
Stable water isotopes ($\delta^{18}$O, $\delta^2$H) follow the Local Meteoric Water Line (LMWL): $\delta^2\text{H} = a + b \cdot \delta^{18}\text{O}$. Deuterium excess $d = \delta^2\text{H} - 8 \cdot \delta^{18}\text{O}$ quantifies deviation.

\subsubsection{Isotope-Based Penalties}
For evaporation, we expect $|E_v| > |E_u|$ (deviation from LMWL) and $d_v < d_u$ (decrease in d-excess). Penalties are added to the transport optimization to bias model selection.

\subsubsection{Nitrate Source Discrimination}
We use two methods for nitrate source attribution:
\begin{enumerate}
	\item \textbf{Dual Isotope Mixing}: Bayesian classification using $\delta^{15}$N and $\delta^{18}$O of nitrate against source endmembers (Manure, Fertilizer, Soil N, Atmospheric).
	\item \textbf{Hydrochemical Ratios}: When isotopes are unavailable, Compositional Data Analysis (CoDA) and ratios ($NO_3/Cl$, $NO_3/K$) are used with Bayesian evidence accumulation.
\end{enumerate}

\subsection{Vadose Zone Physics}\label{sec:vadose_physics}
The optional vadose-zone module simulates 1D unsaturated flow using Richards' equation:
\begin{equation}\label{eq:richards}
	C(\psi) \frac{\partial \psi}{\partial t} = \frac{\partial}{\partial z} \left[ K(\psi) \left( \frac{\partial \psi}{\partial z} - 1 \right) \right] - S(\psi,z,t),
\end{equation}
where $C(\psi) = \partial \theta / \partial \psi$ is the specific moisture capacity.

\subsubsection{Numerical Discretization}
We discretize \eqref{eq:richards} using an implicit finite difference scheme on a cell-centered grid with spacing $\Delta z$. Using the backward Euler approximation for time:
\begin{equation}
	C_i^{n+1, m} \frac{\psi_i^{n+1, m+1} - \psi_i^n}{\Delta t} = \frac{1}{\Delta z} \left[ q_{i-1/2}^{n+1, m} - q_{i+1/2}^{n+1, m} \right] - S_i
\end{equation}
where $m$ is the Picard iteration index. The inter-node Darcy fluxes $q$ are computed using harmonic means for hydraulic conductivity $K$:
\begin{equation}
	q_{i+1/2} = -K_{i+1/2} \left[ \frac{\psi_{i+1} - \psi_i}{\Delta z} - 1 \right], \quad K_{i+1/2} = \frac{2 K_i K_{i+1}}{K_i + K_{i+1}}
\end{equation}
The numerical implementation employs a modified Picard iteration with robust boundary condition handling. For Dirichlet (constant head) bottom boundaries, the Darcy flux is computed as $q = K(1 - \frac{\psi_{bc} - \psi_{prev}}{\Delta z})$, while free drainage assumes a unit gradient ($q = K$).

Outputs include vadose-to-groundwater travel-time priors and nitrate breakthrough curves via dual-domain TTD convolution.

\subsubsection{Quasi-2D Extensions (Tensor Anisotropy and Seepage Face)}
To approximate flow in dipping stratigraphy without a full 2D mesh, Hydrosheaf implements a tensor projection for vertical conductivity. For a layer with dip angle $\alpha$ and anisotropy ratio $A = K_{\parallel}/K_{\perp}$:
\begin{equation}
	K_{zz}(\psi) = K_{\perp}(\psi) \left( A \sin^2 \alpha + \cos^2 \alpha \right)
\end{equation}
This effectively scales the saturated conductivity $K_s$ to represent the vertical component of the hydraulic tensor, allowing 1D columns to simulate infiltration in sloping aquifers.

Additionally, a dynamic \textbf{Seepage Face} boundary condition is implemented to model discharge to surface drains. The boundary switches between a Neumann condition ($q=0$) when unsaturated ($h < 0$) and a Dirichlet condition ($h=0$) when saturated, iteratively updating the status based on computed flux direction within the Picard loop.

\subsubsection{Transverse Dispersion via Lateral Topology}
Standard graph models restrict transport to 1D advection along edges, ignoring transverse dispersion ($D_T \nabla^2 C$) which spreads mass perpendicular to the flow direction. To address this without a mesh, Hydrosheaf employs a \textbf{Lateral Topology} inference step.

Neighbors that are spatially close ($d < r$) but hydraulically equipotential ($\nabla h \approx 0$) are classified as ``Lateral Edges''. These edges do not carry advective flux but serve as potential mixing sources. For a primary flow path $u \to v$, the solver evaluates mixing models where the downstream composition $x_v$ is a mixture of the upstream water $x_u$ and any lateral neighbors $x_{lat}$:
\begin{equation}
	x_v = (1-f) x_u + f x_{lat} + \Delta x_{rxn}
\end{equation}
This effectively discretizes the transverse dispersion term into a finite-volume mixing interaction, ensuring that dispersion is only inferred when supported by observed lateral gradients.

\newpage

% SECTION 3
\section{Numerical Solution Strategy}\label{sec:numerical_solution}

\subsection{Network Optimization Algorithm}\label{sec:network_optimization}
The core algorithm employs a nested evaluation strategy. For each edge and each transport model candidate $c \in \mathcal{C}$:
\begin{enumerate}
	\item \textbf{Transport Optimization}: Compute $\boldsymbol{\theta}_c^*$ using conservative weights.
	\item \textbf{Reaction Optimization}: Solve the bounded LASSO problem for $\mathbf{z}_c^*$ using full weights.
	\item \textbf{Scoring}: Compute total objective $J_c$.
\end{enumerate}
The best model is selected via Boltzmann-weighted probabilities: $p_c \propto \exp(-J_c)$.

\subsection{Solution of Reaction Equations}\label{sec:solution_reaction}
Problem \eqref{eq:lasso} is solved via Coordinate Descent (Algorithm \ref{alg:coord_descent}). To ensure numerical stability, the algorithm includes an adaptive ridge regularization parameter $\epsilon \approx 10^{-10}$ and performs a multicollinearity check.

Specifically, we estimate the condition number $\kappa$ of the Gram matrix $G = S^\top \mathbf{W} S$. If $\kappa > 30$, a warning is issued indicating that the reaction set may be ill-conditioned, likely due to linear dependence between mineral stoichiometries (e.g., competing silicate weathering proxies).

\begin{algorithm}
	\caption{Coordinate Descent for Bounded LASSO}\label{alg:coord_descent}
	\begin{algorithmic}[1]
		\State \textbf{Input:} Residual $\mathbf{r}$, stoichiometric matrix $S$, weights $\mathbf{W}$, penalty $\lambda$, bounds $\boldsymbol{\ell}, \mathbf{u}$
		\Repeat
		\For{$j = 1$ to $m$}
		\State Compute partial residual $\rho_j$ and normalization $\nu_j$
		\State Update $z_j \leftarrow \max\{\ell_j, \min\{\mathcal{S}(\rho_j, \lambda/2) / \nu_j, u_j\}\}$
		\EndFor
		\Until{convergence}
	\end{algorithmic}
\end{algorithm}

\begin{theorem}[Convergence of Coordinate Descent]\label{thm:cd_convergence}
	For the strictly convex bounded LASSO problem in \eqref{eq:lasso}, the coordinate descent algorithm (Algorithm \ref{alg:coord_descent}) converges to the unique global minimizer.
\end{theorem}

\begin{proof}
	The objective function is strictly convex due to the positive definite weight matrix $\mathbf{W}$, and the feasible set is compact. Each coordinate update decreases the objective (or leaves it unchanged), ensuring monotonic convergence. Global convergence follows from standard results on block coordinate descent for convex problems.
\end{proof}

Stability-based tuning selects $\lambda$ via resampling.

\subsection{Graph Inference Algorithms}\label{sec:graph_algorithms}
\subsubsection{Bayesian Head Estimation}
When direct head measurements are missing, a Bayesian linear-Gaussian model estimates hydraulic heads from topography and depth-to-water priors, propagating uncertainty to edge probabilities.

\subsubsection{Sheaf-Consistent Edge Refinement}
A sheaf-consistency score $J^{\text{sheaf}}_{uv}$ blends head priors, isotope compatibility, and conservative tracer behavior. The evaporation probability $\pi_{uv}$ is refined using three physical heuristics:
\begin{enumerate}
	\item \textbf{Depth Constraint}: Evaporation is attenuated with depth ($z > 30$m).
	\item \textbf{Thermodynamic Directionality}: $\pi_{uv}$ is reduced if $|E_v| \leq |E_u|$, as evaporation must increase the evaporation index.
	\item \textbf{d-Excess Signature}: Kinetic fractionation uniquely decreases $d$-excess; paths showing $d$-excess increases are penalized as evaporation candidates.
\end{enumerate}
Soft selection (Expectation-Maximization) iteratively updates edge weights to ensure global consistency.

\subsection{Temporal Deconvolution}\label{sec:temporal_deconvolution}
Time-varying signals are resolved via:
\begin{itemize}
	\item \textbf{Cross-Correlation}: Estimates mean residence time $\tau^*$ from tracer time series.
	\item \textbf{TTD Convolution}: Models downstream concentration as
	      \[C_v(t) \approx \int C_u(t-\tau) g(\tau) e^{-k\tau} d\tau,\]
	      solving for the kernel $g(\tau)$ via nonnegative least squares. For the Piston Flow Model (PFM), where $g(\tau) = \delta(\tau - \tau^*)$, the implementation includes physical validity checks to warn for negative residence times or extrapolation beyond the available input history.
\end{itemize}

\newpage

% SECTION 4
\section{Inverse Modeling and Parameter Estimation}\label{sec:calibration}

\subsection{Definition of the Objective Function}\label{sec:objective_function}
For each edge $(u,v)$, we minimize:
\begin{equation}\label{eq:master_problem}
	\mathcal{L}(\boldsymbol{\theta}, \mathbf{z}) = \|\mathbf{x}_v - \mathbf{x}_{\text{pred}}\|_{\mathbf{W}}^2 + \lambda \|\mathbf{z}\|_1 + \mathcal{P}_{\text{EC/TDS}} + \mathcal{P}_{\text{iso}} + \mathcal{P}_{\text{Gibbs}} + \mathcal{P}_{\text{cons}}
\end{equation}
subject to bounds $\boldsymbol{\ell} \leq \mathbf{z} \leq \mathbf{u}$ and $\boldsymbol{\theta} \in \Theta$.

\subsection{Optimization Algorithms}\label{sec:optimization_algorithms}
For system-wide calibration (e.g. recharge, decay rates), Hydrosheaf supports:
\begin{itemize}
	\item \textbf{Native GLM Engine}: Gauss-Levenberg-Marquardt with trust-region reflective approach.
	\item \textbf{PEST++ Integration}: Wrapper for PEST++-GLM and PEST++-IES for highly parameterized inversion.
\end{itemize}

\subsection{Uncertainty Quantification}\label{sec:uncertainty}
\subsubsection{Bayesian MCMC Formulation}
We formulate the inverse problem in a probabilistic framework to capture the joint posterior distribution of transport and reaction parameters. The model is:
\begin{equation}
	\mathbf{x}_v = \gamma \mathbf{x}_u + S \mathbf{z} + \boldsymbol{\epsilon}, \quad \boldsymbol{\epsilon} \sim \mathcal{N}(\mathbf{0}, \Sigma)
\end{equation}
where $\Sigma = \text{diag}(\sigma_1^2, \dots, \sigma_n^2)$ represents heteroscedastic measurement and model structural error. The joint posterior is:
\begin{equation}
	P(\gamma, \mathbf{z}, \boldsymbol{\sigma} | \mathbf{x}_u, \mathbf{x}_v) \propto P(\mathbf{x}_v | \gamma, \mathbf{z}, \boldsymbol{\sigma}) P(\gamma) P(\mathbf{z}) P(\boldsymbol{\sigma})
\end{equation}
We employ a \textbf{Laplace Prior} for the reaction extents to enforce sparsity, consistent with the L1 penalty in the deterministic LASSO: $P(z_j) = \frac{1}{2b} \exp(-|z_j|/b)$.

The sampler uses the No-U-Turn Sampler (NUTS) via PyMC. Thermodynamic bounds are implemented using \texttt{pm.Truncated} distributions rather than hard potential penalties to maintain smooth gradients for the HMC sampler.

\subsubsection{Residual Bootstrap}
Bias-Corrected Accelerated (BCa) bootstrap resamples residuals to generate confidence intervals. To ensure physical consistency across replicates, the transport model (e.g., specific endmember or evaporation) is fixed to the original best-fit selection during resampling, preventing the mixing of dissimilar parameter distributions.


\newpage

% SECTION 5
\section{Problem Definition and Program Implementation}\label{sec:implementation}

\subsection{Implementation Map}\label{sec:impl_map}
The implementation is modular. Key mappings:
\begin{itemize}
	\item Edge Inference: \texttt{hydrosheaf/graph/build.py}
	\item Transport/Reaction Fitting: \texttt{hydrosheaf/inference/edge\_fit.py}
	\item Thermodynamics: \texttt{hydrosheaf/phreeqc/}
	\item Nitrate Source: \texttt{hydrosheaf/nitrate\_source\_v2.py}
\end{itemize}

\subsection{Programmatic API}\label{sec:api}
Primary entry points:
\begin{verbatim}
from hydrosheaf.inference.network_fit import infer_edges, fit_network
edges = infer_edges(samples, method="probabilistic_map", config=config)
results = fit_network(samples, edges, config=config)
\end{verbatim}

\subsection{Computational Complexity}\label{sec:complexity}
Total complexity is $O(|E| \cdot (|\mathcal{C}| \cdot n + T \cdot m \cdot n))$. For a medium network (120 edges), runtime is $\approx 15$ min on standard hardware.

\subsection{Extension Capabilities}\label{sec:extensions}
\begin{itemize}
	\item \textbf{Reactive Transport}: Kinetic validation via Damköhler number check.
	\item \textbf{3D Networks}: Layer-aware mineralogy and anisotropic distance.
	\item \textbf{Physics Priors}: Integration with MODPATH endpoints.
\end{itemize}

\newpage

% SECTION 6
\section{Example Problems and Validation}\label{sec:results}

\subsection{Validation Methodology}\label{sec:validation}
Validation criteria include: Charge Balance Error (<5\%), EC/TDS consistency, Cross-validation ($R^2$), Thermodynamic consistency, and Isotope coherence.

\subsection{Representative Case Studies}\label{sec:case_studies}
\subsubsection{Salinization in Irrigated Aquifers}
Differentiation between evaporation ($\gamma > 1$, isotope enrichment) and halite dissolution (stoichiometric Na/Cl release).

\subsubsection{Nitrate Contamination and Denitrification}
Identification of denitrification via negative $NO_3$ extent coupled with $HCO_3$ production.

\subsection{Benchmarks}\label{sec:benchmarks}
Performance scales linearly with edge count. LASSO solver dominates cost but is efficient.

\newpage

% SECTION 7
\section{Input Data Reference}\label{sec:input_data}

\subsection{Command Line Interface}\label{sec:cli}
\begin{itemize}
	\item \texttt{--samples}: Path to samples CSV.
	\item \texttt{--infer-edges}: Enable topology inference.
	\item \texttt{--temporal-enabled}: Enable time-series analysis.
	\item \texttt{--nitrate-source-enabled}: Enable source discrimination.
	\item \texttt{--3d}: Enable 3D mode.
\end{itemize}

\subsection{Data Schemas}\label{sec:schemas}
\begin{itemize}
	\item \textbf{Samples CSV}: Must contain \texttt{id} and ion columns (Ca, Mg, ...). Optional: \texttt{head}, \texttt{elevation}, isotopes.
	\item \textbf{Temporal CSV}: Long-form with \texttt{site\_id}, \texttt{timestamp}, \texttt{parameter}, \texttt{value}.
\end{itemize}

\newpage

% SECTION 8
\section{Output Data Reference}\label{sec:output_data}

\subsection{Output Files}\label{sec:output_files}
\begin{itemize}
	\item JSON or CSV format.
	\item One row per directed edge $(u \to v)$.
\end{itemize}

\subsection{Result Tables Structure}\label{sec:result_tables}
Columns include:
\begin{itemize}
	\item \texttt{edge\_id}, \texttt{u}, \texttt{v}: Topology.
	\item \texttt{transport\_model}, \texttt{gamma}, \texttt{f}: Transport parameters.
	\item \texttt{reaction\_*}: Reaction extents (mmol/L).
	\item \texttt{chemistry\_r2}, \texttt{objective\_score}: Fit metrics.
	\item \texttt{qc\_flags}, \texttt{phreeqc\_ok}: Diagnostics.
\end{itemize}

\newpage

% REFERENCES
% This is a new section to be added to hydrosheaf_technical_document.tex before the Conclusion section

\newpage
\section{Web-Based User Interface}\label{sec:webui}

To make Hydrosheaf accessible to a broader audience beyond command-line users, we have developed a modern full-stack web application (version 0.4.0). This interface provides interactive visualization, project management, and real-time analysis capabilities while maintaining the mathematical rigor of the underlying computational framework.

\subsection{Architecture Overview}

The web application follows a three-tier architecture:

\begin{itemize}
	\item \textbf{Frontend Layer}: React-based single-page application (SPA) with modern UI components, state management via React Context, and real-time updates via WebSockets.
	\item \textbf{Backend Layer}: FastAPI (Python 3.8+) REST API providing endpoints for project management, sample data validation, edge inference, and asynchronous analysis execution.
	\item \textbf{Computational Core}: The same mathematical framework presented in Sections \ref{sec:math_framework}--\ref{sec:extensions}, ensuring consistency between CLI and web interfaces. All transport optimization, LASSO solving, and thermodynamic constraints use identical implementations.
\end{itemize}

\subsection{Frontend Features}

The React-based frontend implements a "Dark Ocean" design theme with glassmorphism effects, providing an intuitive user experience for complex hydrogeochemical analysis:

\subsubsection{Interactive Dashboard}

The main dashboard provides:
\begin{itemize}
	\item \textbf{Project Management}: Create, load, and manage multiple groundwater analysis projects with persistent browser storage (IndexedDB) or server-side PostgreSQL backend.
	\item \textbf{Quick Statistics}: Real-time summaries of sample counts, inferred edges, and analysis completion status.
	\item \textbf{Recent Activity}: Timeline view of analysis runs, parameter changes, and results exports with version history.
\end{itemize}

\subsubsection{Sample Data Manager}

Provides spreadsheet-like interface for:
\begin{itemize}
	\item \textbf{Data Entry}: Input or paste chemical concentrations, coordinates, and metadata with automatic unit conversion (mg/L $\leftrightarrow$ mmol/L, meq/L $\leftrightarrow$ mmol/L).
	\item \textbf{CSV Import/Export}: Bulk upload of field data with schema validation (via Pydantic models) and error highlighting.
	\item \textbf{Quality Control}: Automated charge balance error calculation, outlier detection via robust Z-scores, and completeness validation (missing Ca, Mg, etc.).
\end{itemize}

\subsubsection{Network Visualization}

Interactive graph editor using D3.js force-directed layouts:
\begin{itemize}
	\item \textbf{Node Positioning}: Drag-and-drop arrangement with optional geographic coordinate alignment (lon/lat $\to$ Web Mercator projection).
	\item \textbf{Edge Inference Controls}: Visual configuration of probabilistic edge inference parameters (radius km, minimum head gradient, $p_{min}$ threshold).
	\item \textbf{3D Layer Support}: Toggle vertical stratification view and visualize cross-layer connectivity (requires \texttt{--3d} mode enabled).
	\item \textbf{Edge Attributes}: Color-coded edges by confidence probability, dominant transport model (evaporation/mixing), or total reaction extent $\|\mathbf{z}\|_1$.
\end{itemize}

\subsubsection{Analysis Configuration Panel}

Formalized interface for model parameters mapped to CLI flags (Table \ref{tab:impl_map}):
\begin{itemize}
	\item \textbf{Transport Models}: Select candidate endmembers (rainfall, seawater, river), enable/disable evaporation hypothesis, configure isotope penalty weights ($\eta_E$, $\eta_d$).
	\item \textbf{Reaction Dictionary}: Toggle minerals (calcite, gypsum, halite, silicates), set custom stoichiometry, select PHREEQC thermodynamic database (WATEQ4F, PHREEQC.dat, llnl.dat).
	\item \textbf{Regularization}: Adjust L1 penalty ($\lambda$), enable residence-time coupling (\texttt{--residence-coupling}), set coordinate descent convergence tolerance.
	\item \textbf{Advanced Options}: Enable temporal mode (\texttt{--temporal-enabled}), select uncertainty quantification method (Bootstrap/Bayesian MCMC), upload vadose priors CSV.
\end{itemize}

\subsubsection{Real-Time Analysis Monitor}

Asynchronous execution with live feedback via WebSockets (FastAPI's \texttt{WebSocketEndpoint}):
\begin{itemize}
	\item \textbf{Progress Tracking}: Per-edge status (queued, running, complete, failed) with estimated time remaining based on historical per-edge solve times.
	\item \textbf{Intermediate Results}: Streaming updates of transport model selection, LASSO iteration count, reaction convergence flags, and current objective value.
	\item \textbf{Log Viewer}: Real-time display of solver diagnostics, PHREEQC warnings (saturation index failures, database issues), and QC flags (charge balance errors, isotope inconsistencies).
\end{itemize}

\subsubsection{Results Dashboard}

Comprehensive visualization of analysis outputs:
\begin{itemize}
	\item \textbf{Summary Tables}: Sortable, filterable grids (AG Grid React) of edge results with dynamic columns: transport parameters ($\gamma$, $f$, endmember ID), reaction extents (one column per active mineral), fit metrics ($R^2$, residual norm, L1 norm).
	\item \textbf{Spatial Maps}: Leaflet-based choropleth visualization of transport model prevalence, TDS evolution gradients, nitrate source probabilities ($P_{\text{manure}}$).
	\item \textbf{Piper/Stiff Diagrams}: Automated hydrochemical facies classification using ternary projections and evolution trajectories along flow paths.
	\item \textbf{Time-Series Plots}: When temporal mode is enabled (\texttt{--temporal-enabled}), interactive Plotly charts of residence time estimates, seasonal decomposition (trend + periodic + residual), and TTD kernel visualizations (histogram of $g(\tau)$).
	\item \textbf{Uncertainty Bands}: Credible intervals (Bayesian MCMC) or bootstrap percentiles (BCa) overlaid on predicted concentrations when \texttt{--uncertainty} is enabled.
\end{itemize}

\subsection{Backend API}

The FastAPI backend exposes RESTful endpoints and WebSocket channels designed around resource-oriented patterns:

\subsubsection{Project Management Endpoints}

\begin{verbatim}
POST   /api/projects              Create new project
GET    /api/projects/{id}         Retrieve project metadata
PUT    /api/projects/{id}         Update configuration
DELETE /api/projects/{id}         Delete project
GET    /api/projects/{id}/export  Export as PDF report
\end{verbatim}

\subsubsection{Sample Data Endpoints}

\begin{verbatim}
POST /api/samples           Upload CSV or JSON sample data
GET  /api/samples/validate  Run schema and QC checks
PUT  /api/samples/{id}      Update individual sample
\end{verbatim}

\subsubsection{Security Architecture}

The backend implements several security measures to protect the integrity of scientific data and computational resources:

\begin{itemize}
	\item \textbf{Rate Limiting}: All API endpoints are protected by a token-bucket rate limiter (default 200 requests/minute) using \texttt{slowapi} to prevent abuse and denial-of-service attacks.
	\item \textbf{Input Sanitization}: Project identifiers and file paths are rigorously sanitized to prevent directory traversal attacks.
	\item \textbf{Validation Layer}: Pydantic schemas enforce strict type checking and range validation on all inputs before they reach the core computational engine.
	\item \textbf{Error Handling}: Production error handlers suppress stack traces to prevent information leakage while logging detailed diagnostics for administrators.
\end{itemize}

\subsubsection{Network Endpoints}


\begin{verbatim}
POST /api/edges/infer  Trigger probabilistic edge inference
GET  /api/edges         Retrieve inferred edges with p(i→j)
PUT  /api/edges/{id}    Manually override edge (enable/disable)
\end{verbatim}

\subsubsection{Analysis Endpoints}

\begin{verbatim}
POST /api/analysis/run            Start analysis (returns job_id)
WS   /api/analysis/ws/{job_id}   WebSocket for real-time progress
GET  /api/analysis/status/{job_id}  Poll job status
GET  /api/analysis/results/{job_id} Retrieve completed results
\end{verbatim}

\subsubsection{Data Transformation Layer}

The backend implements automatic conversions to maintain compatibility with the core framework:
\begin{itemize}
	\item \textbf{Unit Conversion}: Frontend accepts mg/L (common in water quality reports), backend converts to mmol/L using molar masses (e.g., $M_{\text{Ca}} = 40.08$ g/mol) before invoking \texttt{hydrosheaf.inference.network\_fit}.
	\item \textbf{Schema Validation}: Pydantic models enforce required fields (\texttt{site\_id}, \texttt{Ca}, \texttt{Mg}, \texttt{Na}, \texttt{HCO3}, \texttt{Cl}, \texttt{SO4}, \texttt{EC}, \texttt{TDS}, \texttt{pH}) and reject malformed inputs with HTTP 422 before reaching the solver.
	\item \textbf{Result Serialization}: Converts NumPy arrays (\texttt{ndarray}) and NetworkX graphs (\texttt{DiGraph}) to JSON-serializable formats (lists of dicts) for frontend consumption.
\end{itemize}

\subsection{Deployment and Scalability}

\subsubsection{Development Mode}

For local testing and demonstration:
\begin{verbatim}
# Terminal 1 (Backend)
cd web/backend
uvicorn app.main:app --reload --port 8000

# Terminal 2 (Frontend)
cd web/frontend
npm run dev  # Vite dev server on port 5173
\end{verbatim}

Open browser to \texttt{http://localhost:5173}. Frontend proxies API requests to \texttt{http://localhost:8000}.

\subsubsection{Production Deployment}

Recommended stack for institutional deployment (university research group, government agency):
\begin{itemize}
	\item \textbf{Reverse Proxy}: Nginx or Traefik for HTTPS termination (Let's Encrypt certificates), load balancing, and static asset serving.
	\item \textbf{Application Server}: Gunicorn with Uvicorn workers (4-8 workers for CPU-bound LASSO solving) for concurrent request handling.
	\item \textbf{Task Queue}: Celery + Redis for background LASSO/MCMC jobs exceeding HTTP timeout limits (default 30s for web APIs).
	\item \textbf{Database}: PostgreSQL 12+ for persistent project storage, user sessions, and analysis result caching. SQLite acceptable for single-user deployments.
	\item \textbf{File Storage}: S3-compatible object storage (MinIO, AWS S3) for large result exports (CSV/JSON $>$ 10 MB) and time-series data.
\end{itemize}

\subsubsection{Performance Considerations}

For networks with $> 100$ edges or Bayesian MCMC uncertainty quantification (4000+ samples per edge):
\begin{itemize}
	\item \textbf{Asynchronous Execution}: Long-running analyses (estimated time $> 30$s) execute in Celery background workers, preventing frontend timeout.
	\item \textbf{Incremental Results}: WebSocket streaming allows users to view partial results (e.g., first 10 edges completed) while remaining edges compute in parallel.
	\item \textbf{Caching}: Redis-backed result cache (TTL 1 hour) prevents redundant computation when users toggle visualization options (Piper vs. Stiff diagrams) without changing model parameters.
	\item \textbf{Database Indexing}: PostgreSQL B-tree indexes on \texttt{project\_id}, \texttt{job\_id}, and \texttt{edge\_id} enable sub-millisecond lookups for result retrieval.
\end{itemize}

\subsection{Integration with Existing Workflows}

The web interface is designed to complement, not replace, the CLI:

\begin{itemize}
	\item \textbf{Bidirectional Compatibility}: Projects created via web UI can be exported as CLI-compatible JSON configuration files (\texttt{hydrosheaf --config exported\_config.json}), and vice versa (CLI results can be imported into web UI for visualization).
	\item \textbf{Programmatic API}: Advanced users can bypass the web frontend and interact directly with the FastAPI backend via Python \texttt{requests} library or R \texttt{httr} package for automated workflows.
	\item \textbf{Batch Processing}: For large regional studies (hundreds of wells, dozens of time steps, MCMC uncertainty), the CLI remains preferred due to lower overhead and easier integration with HPC schedulers (SLURM, PBS).
\end{itemize}

\subsection{Accessibility and Usability}

The web interface follows modern accessibility best practices:
\begin{itemize}
	\item \textbf{WCAG 2.1 Level AA Compliance}: Sufficient color contrast ratios (4.5:1 for text), keyboard navigation support, ARIA labels for screen readers.
	\item \textbf{Responsive Design}: Mobile-friendly layouts (Tailwind CSS breakpoints) for tablet and phone access, though desktop recommended for complex network editing.
	\item \textbf{Internationalization (i18n)}: React-i18next framework supports multiple languages (English, Spanish, French planned for future releases).
\end{itemize}

\subsection{Future Development}

Planned enhancements to the web interface (see \texttt{web/INTEGRATION\_ROADMAP.md} for detailed timeline):

\begin{itemize}
	\item \textbf{Collaborative Features}: Multi-user access with role-based permissions (viewer, analyst, admin) using OAuth2/OIDC authentication (Keycloak, Auth0).
	\item \textbf{Interactive Tutorials}: Guided walkthroughs using synthetic datasets (similar to \texttt{hydrosheaf\_synthetic\_csv/}) to teach inverse modeling concepts interactively.
	\item \textbf{3D Visualization}: WebGL-based rendering (Three.js) of layered aquifer networks and vertical transport pathways, replacing 2D projections.
	\item \textbf{Parameter Sensitivity Analysis}: Interactive dashboards showing how $\lambda$, weights $\mathbf{W}$, and thermodynamic bounds affect sparsity $|\mathcal{A}|$ and fit quality $R^2$.
	\item \textbf{Integration with External Databases}: Connectors for USGS Water Quality Portal API, state monitoring databases (e.g., California GAMA), and remote PHREEQC thermodynamic data repositories.
	\item \textbf{Live Model Comparison}: Side-by-side comparison of multiple analysis runs (e.g., varying $\lambda$ from 0.01 to 1.0) with synchronized visualizations.
\end{itemize}

\subsection{Open-Source Licensing and Contributions}

The web application (both frontend and backend) is released under the same MIT License as the core Hydrosheaf package. Contributions welcome via GitHub pull requests. Frontend uses npm for dependency management (React 18, Vite 4, Tailwind CSS 3); backend uses Poetry for Python dependency management (FastAPI 0.104, Pydantic 2.5).



\section*{Acknowledgments}
This work synthesizes methods from convex optimization, geochemistry, and hydrogeology.

\begin{thebibliography}{99}
	\bibitem{lasaga1998} Lasaga, A. C. (1998). \textit{Kinetic Theory in the Earth Sciences}.
	\bibitem{gelman2013} Gelman, A., et al. (2013). \textit{Bayesian Data Analysis}.
	\bibitem{efron1987} Efron, B. (1987). Better Bootstrap Confidence Intervals.
	\bibitem{friedman2007} Friedman, J., et al. (2007). Pathwise coordinate optimization.
	\bibitem{friedman2010} Friedman, J., et al. (2010). Regularization paths for generalized linear models.
	\bibitem{parkhurst2013} Parkhurst, D. L., \& Appelo, C. A. J. (2013). PHREEQC version 3.
	\bibitem{plummer1994} Plummer, L. N., et al. (1994). NETPATH version 2.0.
	\bibitem{gwb} Aqueous Solutions LLC. (n.d.). \textit{The Geochemist's Workbench}.
	\bibitem{wang2025} Wang, Y., et al. (2025). Enhancing inverse modeling in groundwater systems.
	\bibitem{gholami2024} Gholami, R., et al. (2024). Geochemistry and machine learning.
	\bibitem{harbaugh2005} Harbaugh, A. W. (2005). MODFLOW-2005.
	\bibitem{pollock2016} Pollock, D. W. (2016). MODPATH Version 7.
	\bibitem{bakker2016} Bakker, M., et al. (2016). Scripting MODFLOW using FloPy.
	\bibitem{pymc2023} PyMC Developers. (2023). PyMC.
	\bibitem{arviz2023} ArviZ Developers. (2023). ArviZ.
	\bibitem{vangenuchten1980} van Genuchten, M. T. (1980). Hydraulic conductivity of unsaturated soils.
	\bibitem{mualem1976} Mualem, Y. (1976). Predicting hydraulic conductivity.
	\bibitem{feddes1978} Feddes, R. A., et al. (1978). Simulation of Field Water Use.
	\bibitem{richards1931} Richards, L. A. (1931). Capillary conduction.
	\bibitem{robinson2008} Robinson, M. (2008). \textit{Sheaves in Geometry and Logic}.
	\bibitem{appelo2005} Appelo, C. A. J., \& Postma, D. (2005). \textit{Geochemistry, Groundwater and Pollution}.
	\bibitem{clark2015} Clark, I. D., \& Fritz, P. (2015). \textit{Environmental Isotopes in Hydrogeology}.
	\bibitem{kendall1998} Kendall, C. (1998). Tracing nitrogen sources.
	\bibitem{xue2009} Xue, D., et al. (2009). Application of stable isotopes.
	\bibitem{curry2014} Curry, J. M. (2014). Sheaves, cosheaves and applications.
	\bibitem{tseng2001} Tseng, P. (2001). Convergence of coordinate descent.
\end{thebibliography}

\end{document}
